\expandafter\gdef\csname SNside1\endcsname{\node[anchor=north west,inner sep=0pt,text width=48mm,font=\SideFont,align=left] at ([xshift=138mm,yshift=-466.85pt]current page.north west) {\textcolor{SNink}{\textbf{最小反例 · 地址宽与数据宽可不同}}\par\vspace{1.5mm}若按字节编址、可寻址64MiB，地址需要26位；若一次访存交换32位，MDR可为32位。MAR为26位而MDR为32位完全自洽，不能因二者都接主存就要求等宽。};}
\expandafter\gdef\csname SNside2\endcsname{\node[anchor=north west,inner sep=0pt,text width=48mm,font=\SideFont,align=left] at ([xshift=138mm,yshift=-459.99pt]current page.north west) {\textcolor{SNink}{\textbf{错项排除 · CPI同为1也不等速}}\par\vspace{1.5mm}单周期CPU一拍必须容纳最慢完整指令，流水线一拍只需容纳最慢段及寄存器开销。二者CPI都可约为1，时钟周期却不同，执行时间还须乘指令数与每拍时长。};}
\expandafter\gdef\csname SNside3\endcsname{\node[anchor=north west,inner sep=0pt,text width=48mm,font=\SideFont,align=left] at ([xshift=138mm,yshift=-379.63pt]current page.north west) {\textcolor{SNink}{\textbf{不变量 · 总线同拍只有一个驱动源}}\par\vspace{1.5mm}单内部总线上，R1送Y与R2送MAR要求两个不同值同时驱动，发生冲突；R1同时送Y和MAR则可能允许，前提是两目标都具备写入路径和独立写使能。};}
\expandafter\gdef\csname SNside4\endcsname{\node[anchor=north west,inner sep=0pt,text width=48mm,font=\SideFont,align=left] at ([xshift=138mm,yshift=-95.00pt]current page.north west) {\textcolor{SNink}{\textbf{最小反例 · 同一动作不可重复计延迟}}\par\vspace{1.5mm}若题给从Read到主存数据送MDR共5拍，这5拍已经含送入MDR。再加一次``写MDR''会重复计数；但随后MDR经内部总线写IR仍是另一项真实传送。};}
\expandafter\gdef\csname SNside5\endcsname{\node[anchor=north west,inner sep=0pt,text width=48mm,font=\SideFont,align=left] at ([xshift=138mm,yshift=-95.00pt]current page.north west) {\textcolor{SNink}{\textbf{错项排除 · 编码对象先分清}}\par\vspace{1.5mm}32条机器指令只需5位区分操作码，不说明控制存储器地址也只需5位。公共2步加每指令4步共有130位置，微地址需要8位；操作码、微地址和控制位数分别统计。};}
\expandafter\gdef\csname SNside6\endcsname{\node[anchor=north west,inner sep=0pt,text width=48mm,font=\SideFont,align=left] at ([xshift=138mm,yshift=-95.00pt]current page.north west) {\textcolor{SNink}{\textbf{最小反例 · 4条命令不只4种状态}}\par\vspace{1.5mm}一个互斥字段含4个微命令，仍需``不发该组命令''的第5种状态，所以通常至少3位。互斥只限制同字段同时发一条，不能省掉本拍暂不用该字段的状态。};}
\expandafter\gdef\csname SNside7\endcsname{\node[anchor=north west,inner sep=0pt,text width=48mm,font=\SideFont,align=left] at ([xshift=138mm,yshift=-95.00pt]current page.north west) {\textcolor{SNink}{\textbf{课程连接 · 缺页使当前指令重试}}\par\vspace{1.5mm}OS调入缺页后，原访问尚未完成，处理器通常从该指令重试；设备完成中断对应一条已完成指令的边界。返回点取决于事件与ISA的精确状态规则，不由当前取指PC猜测。};}
\expandafter\gdef\csname SNside8\endcsname{\node[anchor=north west,inner sep=0pt,text width=48mm,font=\SideFont,align=left] at ([xshift=138mm,yshift=-95.00pt]current page.north west) {\textcolor{SNink}{\textbf{错项排除 · 进入内核不等于换进程}}\par\vspace{1.5mm}一次外设中断可能保存现场、处理后仍返回同一用户进程。模式切换、服务程序运行和调度上下文切换分别属于权限、控制流和执行主体三个层面，题目未要求时不合并。};}
\expandafter\gdef\csname SNside9\endcsname{\node[anchor=north west,inner sep=0pt,text width=48mm,font=\SideFont,align=left] at ([xshift=138mm,yshift=-95.00pt]current page.north west) {\textcolor{SNink}{\textbf{不变量 · 第一条付满延迟，后续按间隔}}\par\vspace{1.5mm}理想k段、n条指令总时长为(k+n\ensuremath{-}1)拍：第一条需k拍，剩余n\ensuremath{-}1条每拍完成一条。不能给每条都加一次填充延迟，也不能把阶段数直接当任何基线下的加速比。};}
\expandafter\gdef\csname SNside10\endcsname{\node[anchor=north west,inner sep=0pt,text width=48mm,font=\SideFont,align=left] at ([xshift=138mm,yshift=-95.00pt]current page.north west) {\textcolor{SNink}{\textbf{错项排除 · 段寄存器还要携带控制}}\par\vspace{1.5mm}只传操作数不传目的编号，结果可能写错寄存器；保持数据却让写使能继续推进，可能重复写入。流水段状态要同时说明值、目的、有效位与控制信息，不只数``几个数据寄存器''。};}
\expandafter\gdef\csname SNside11\endcsname{\node[anchor=north west,inner sep=0pt,text width=48mm,font=\SideFont,align=left] at ([xshift=138mm,yshift=-95.00pt]current page.north west) {\textcolor{SNink}{\textbf{错项排除 · 两种“动态”分类角度不同}}\par\vspace{1.5mm}动态流水线讨论同时配置不同功能的能力；动态调度讨论后面已就绪指令能否先执行。某设计可采用其中一项而未采用另一项，不能看到名称相近就把两项机制互相证明。};}
\expandafter\gdef\csname SNside12\endcsname{\node[anchor=north west,inner sep=0pt,text width=48mm,font=\SideFont,align=left] at ([xshift=138mm,yshift=-95.00pt]current page.north west) {\textcolor{SNink}{\textbf{最小反例 · 尚未产生的值不能转发}}\par\vspace{1.5mm}LOAD的值到第5拍MEM末才形成，下一条本应第5拍EX开头使用它，两者时刻不合。已有MEM\ensuremath{\to}EX旁路仍需把消费者推至第6拍；``有转发''不能抹去因果顺序。\par\smallskip\begin{tikzpicture}[x=1mm,y=1mm,>=Stealth,node distance=5mm]
\node[draw=SNink!50,fill=SNblue,rounded corners=1mm,text width=34mm,align=center,inner sep=1.2mm,font=\fontsize{7.5}{9}\selectfont] (m0) at (0,0) {MEM末5};
\node[draw=SNink!50,fill=SNblue,rounded corners=1mm,text width=34mm,align=center,inner sep=1.2mm,font=\fontsize{7.5}{9}\selectfont] (m1) at (0,-11) {EX初5};
\node[draw=SNink!50,fill=SNblue,rounded corners=1mm,text width=34mm,align=center,inner sep=1.2mm,font=\fontsize{7.5}{9}\selectfont] (m2) at (0,-22) {改到EX6};
\draw[->,SNink] (m0) edge[bend left=40] (m2);
\draw[->,SNink] (m1) -- (m2);
\end{tikzpicture}
\par{\fontsize{7.5}{9}\selectfont 同拍初早于同拍末，故需推迟消费者}};}
\expandafter\gdef\csname SNside13\endcsname{\node[anchor=north west,inner sep=0pt,text width=48mm,font=\SideFont,align=left] at ([xshift=138mm,yshift=-95.00pt]current page.north west) {\textcolor{SNink}{\textbf{不变量 · 产生点必须早于使用点}}\par\vspace{1.5mm}先把每个值的产生段末与消费者所需段初标在时间线上，再检查是否有实际转发路径。STORE数据端、分支比较端和ALU输入端可能不同，不能把一种load-use模板用于全部指令。};}
\expandafter\gdef\csname SNside14\endcsname{\node[anchor=north west,inner sep=0pt,text width=48mm,font=\SideFont,align=left] at ([xshift=138mm,yshift=-95.00pt]current page.north west) {\textcolor{SNink}{\textbf{最小反例 · 改名字不能提前算出值}}\par\vspace{1.5mm}前指令产生r1，后指令确实需要这个结果，即使把目的映射为物理寄存器p7，后指令仍须等p7得到值。重命名消除WAR/WAW名字冲突，不能消除真实RAW因果依赖。};}
\expandafter\gdef\csname SNside15\endcsname{\node[anchor=north west,inner sep=0pt,text width=48mm,font=\SideFont,align=left] at ([xshift=138mm,yshift=-95.00pt]current page.north west) {\textcolor{SNink}{\textbf{最小反例 · 一核也能保存两线程状态}}\par\vspace{1.5mm}一个核心可保存两套PC和寄存器状态，按周期或等待事件交替执行，也可SMT跨线程发射。线程状态数与物理核心数不同，硬件多线程不要求先有多个核心。};}
\expandafter\gdef\csname SNside16\endcsname{\node[anchor=north west,inner sep=0pt,text width=48mm,font=\SideFont,align=left] at ([xshift=138mm,yshift=-95.00pt]current page.north west) {\textcolor{SNink}{\textbf{课程连接 · 共享地址并不自动互斥}}\par\vspace{1.5mm}多核共享物理地址空间，仍可能并发读到同一旧值并分别写回。OS同步中的原子操作与锁解决复合更新的协调；Cache一致性不等于所有读改写序列自动原子。};}
\expandafter\gdef\csname SNside17\endcsname{\node[anchor=north west,inner sep=0pt,text width=48mm,font=\SideFont,align=left] at ([xshift=138mm,yshift=-95.00pt]current page.north west) {\textcolor{SNink}{\textbf{不变量 · 调度下界来自不可颠倒链}}\par\vspace{1.5mm}无转发且不许同拍读写时，LOAD\ensuremath{\to}翻倍\ensuremath{\to}相加\ensuremath{\to}STORE链的各次ID必须晚于相应生产者WB。先给关键链下界，再排出达到该下界的合法时序，才能证明``最少几拍''。};}
\expandafter\gdef\csname SNside18\endcsname{\node[anchor=north west,inner sep=0pt,text width=48mm,font=\SideFont,align=left] at ([xshift=138mm,yshift=-95.00pt]current page.north west) {\textcolor{SNink}{\textbf{不变量 · 调度下界来自不可颠倒链}}\par\vspace{1.5mm}无转发且不许同拍读写时，LOAD\ensuremath{\to}翻倍\ensuremath{\to}相加\ensuremath{\to}STORE链的各次ID必须晚于相应生产者WB。先给关键链下界，再排出达到该下界的合法时序，才能证明``最少几拍''。};}
\AddToShipoutPictureFG{\begin{tikzpicture}[remember picture,overlay]\ifcsname SNside\thepage\endcsname\csname SNside\thepage\endcsname\fi\end{tikzpicture}}